\documentclass[11pt, a4paper]{article}
\usepackage[utf8]{inputenc}
\usepackage[margin=0.75in]{geometry}
\usepackage{amsmath, amssymb}
\usepackage{enumitem}
\usepackage{multicol}

\title{\textbf{NEET Chemistry Practice Test: Chemical Kinetics} \\ \large (45 Moderate to Difficult MCQs Covering Complete NCERT Syllabus)}
\date{}

\begin{document}

\maketitle

\section*{Section A: Rate of Reaction, Rate Law, Order \& Molecularity}

\begin{enumerate}

\item For the reaction $2\text{N}_2\text{O}_5(g) \rightarrow 4\text{NO}_2(g) + \text{O}_2(g)$, the rate of formation of $\text{NO}_2$ is $2.8 \times 10^{-3}\text{ M s}^{-1}$. The rate of disappearance of $\text{N}_2\text{O}_5$ and the rate of the reaction respectively are:
\begin{enumerate}
    \item $1.4 \times 10^{-3}\text{ M s}^{-1}$ and $0.7 \times 10^{-3}\text{ M s}^{-1}$
    \item $2.8 \times 10^{-3}\text{ M s}^{-1}$ and $1.4 \times 10^{-3}\text{ M s}^{-1}$
    \item $5.6 \times 10^{-3}\text{ M s}^{-1}$ and $2.8 \times 10^{-3}\text{ M s}^{-1}$
    \item $1.4 \times 10^{-3}\text{ M s}^{-1}$ and $2.8 \times 10^{-3}\text{ M s}^{-1}$
\end{enumerate}

\item In the reaction $\text{BrO}_3^-(aq) + 5\text{Br}^-(aq) + 6\text{H}^+(aq) \rightarrow 3\text{Br}_2(l) + 3\text{H}_2\text{O}(l)$, if the rate of disappearance of $\text{Br}^-$ is $k_1[\text{BrO}_3^-][\text{Br}^-][\text{H}^+]^2$, the rate of appearance of $\text{Br}_2$ is $k_2[\text{BrO}_3^-][\text{Br}^-][\text{H}^+]^2$, and the rate of disappearance of $\text{H}^+$ is $k_3[\text{BrO}_3^-][\text{Br}^-][\text{H}^+]^2$, then the correct relation among $k_1, k_2,$ and $k_3$ is:
\begin{enumerate}
    \item $k_1 = k_2 = k_3$
    \item $\frac{k_1}{5} = \frac{k_2}{3} = \frac{k_3}{6}$
    \item $5k_1 = 3k_2 = 6k_3$
    \item $\frac{k_1}{3} = \frac{k_2}{5} = \frac{k_3}{6}$
\end{enumerate}

\item For a gaseous reaction $A(g) + 2B(g) \rightarrow C(g)$, the rate law is given by $r = k[A][B]^2$. If the volume of the reaction vessel is suddenly reduced to one-fourth of its initial volume, the rate of reaction relative to the initial rate will become:
\begin{enumerate}
    \item $16\text{ times}$
    \item $64\text{ times}$
    \item $\frac{1}{64}\text{ times}$
    \item $32\text{ times}$
\end{enumerate}

\item For the reaction $2A + B \rightarrow C + D$, the following kinetic data were obtained:
\begin{center}
\begin{tabular}{|c|c|c|c|}
\hline
\textbf{Exp.} & $[A]\ (\text{mol L}^{-1})$ & $[B]\ (\text{mol L}^{-1})$ & \textbf{Initial rate of formation of } $D\ (\text{mol L}^{-1}\text{ min}^{-1})$ \\
\hline
I & $0.1$ & $0.1$ & $6.0 \times 10^{-3}$ \\
II & $0.3$ & $0.2$ & $7.2 \times 10^{-2}$ \\
III & $0.3$ & $0.4$ & $2.88 \times 10^{-1}$ \\
IV & $0.4$ & $0.1$ & $2.40 \times 10^{-2}$ \\
\hline
\end{tabular}
\end{center}
The rate law for the reaction is:
\begin{enumerate}
    \item $\text{Rate} = k[A]^2[B]$
    \item $\text{Rate} = k[A][B]^2$
    \item $\text{Rate} = k[A]^2[B]^2$
    \item $\text{Rate} = k[A][B]$
\end{enumerate}

\item The unit of rate constant for a reaction is $\text{atm}^{-3/2}\text{ s}^{-1}$. The order of this reaction is:
\begin{enumerate}
    \item $1.5$
    \item $2.5$
    \item $-0.5$
    \item $0.5$
\end{enumerate}

\item Consider the following statements regarding order and molecularity of a reaction:
\begin{enumerate}[label=(\roman*)]
    \item Order of a reaction can be zero or fractional, whereas molecularity is always a positive integer.
    \item Molecularity is applicable only to elementary reactions, whereas order is applicable to both elementary and complex reactions.
    \item For a complex reaction, the order is always equal to the molecularity of the slowest step.
    \item Reactions with molecularity greater than three are extremely rare.
\end{enumerate}
Which of the above statements are correct?
\begin{enumerate}
    \item (i), (ii) and (iv) only
    \item (i), (ii) and (iii) only
    \item (ii), (iii) and (iv) only
    \item (i), (ii), (iii) and (iv)
\end{enumerate}

\item The decomposition of ozone proceeds via the following two-step mechanism: \\
Step 1: $\text{O}_3(g) \rightleftharpoons \text{O}_2(g) + \text{O}(g)$ \quad (Fast equilibrium) \\
Step 2: $\text{O}_3(g) + \text{O}(g) \rightarrow 2\text{O}_2(g)$ \quad (Slow) \\
The overall order of the reaction with respect to $\text{O}_3$ and $\text{O}_2$ respectively are:
\begin{enumerate}
    \item $2$ and $1$
    \item $2$ and $-1$
    \item $1$ and $-1$
    \item $2$ and $0$
\end{enumerate}

\item For the reaction $2\text{NO}(g) + \text{Cl}_2(g) \rightarrow 2\text{NOCl}(g)$, the mechanism proposed is: \\
Step I: $\text{NO}(g) + \text{Cl}_2(g) \rightleftharpoons \text{NOCl}_2(g)$ \quad (Fast, equilibrium constant $K_c$) \\
Step II: $\text{NOCl}_2(g) + \text{NO}(g) \xrightarrow{k_2} 2\text{NOCl}(g)$ \quad (Slow) \\
The observed rate constant $k_{\text{obs}}$ is equal to:
\begin{enumerate}
    \item $k_2 K_c$
    \item $\frac{k_2}{K_c}$
    \item $k_2 K_c^2$
    \item $\frac{K_c}{k_2}$
\end{enumerate}

\item If the concentration of reactant $A$ is doubled, the rate of reaction increases by a factor of $2\sqrt{2}$. If the concentration of $A$ is quadrupled, the rate of reaction will increase by a factor of:
\begin{enumerate}
    \item $4$
    \item $8$
    \item $16$
    \item $4\sqrt{2}$
\end{enumerate}

\item In the thermal decomposition of acetaldehyde: $\text{CH}_3\text{CHO}(g) \rightarrow \text{CH}_4(g) + \text{CO}(g)$, the rate is proportional to $[\text{CH}_3\text{CHO}]^{3/2}$. If the pressure is expressed in $\text{bar}$ and time in $\text{minutes}$, the units of rate and rate constant respectively are:
\begin{enumerate}
    \item $\text{bar min}^{-1}$ and $\text{bar}^{-1/2}\text{ min}^{-1}$
    \item $\text{bar min}^{-1}$ and $\text{bar}^{1/2}\text{ min}^{-1}$
    \item $\text{bar}^{-1/2}\text{ min}^{-1}$ and $\text{bar min}^{-1}$
    \item $\text{bar min}^{-1}$ and $\text{bar}^{-3/2}\text{ min}^{-1}$
\end{enumerate}

\section*{Section B: Integrated Rate Equations \& Half-Life}

\item Which of the following reactions is an example of a zero-order reaction?
\begin{enumerate}
    \item Acid-catalyzed hydrolysis of ethyl acetate
    \item Inversion of cane sugar
    \item Decomposition of gaseous ammonia on a hot platinum surface at high pressure
    \item Decomposition of $\text{N}_2\text{O}_5$ in $\text{CCl}_4$ solution
\end{enumerate}

\item For a zero-order reaction $A \rightarrow \text{Products}$, the initial concentration of $A$ is $0.2\text{ M}$. If the concentration of $A$ reduces to $0.15\text{ M}$ in $10\text{ minutes}$, the time required for the complete consumption of $A$ is:
\begin{enumerate}
    \item $30\text{ min}$
    \item $40\text{ min}$
    \item $20\text{ min}$
    \item $50\text{ min}$
\end{enumerate}

\item A first-order reaction is $40\%$ complete in $50\text{ minutes}$. In what time will the reaction be $80\%$ complete? ($\log 2 = 0.3010$, $\log 3 = 0.4771$, $\log 5 = 0.6990$)
\begin{enumerate}
    \item $100\text{ minutes}$
    \item $157.5\text{ minutes}$
    \item $125.4\text{ minutes}$
    \item $180.2\text{ minutes}$
\end{enumerate}

\item For a first-order reaction, the ratio of time required for $99.9\%$ completion ($t_{99.9\%}$) to the time required for $90\%$ completion ($t_{90\%}$) is:
\begin{enumerate}
    \item $2$
    \item $3$
    \item $10$
    \item $1.5$
\end{enumerate}

\item The half-life of a reaction is halved when the initial concentration of the reactant is doubled. The order of the reaction is:
\begin{enumerate}
    \item $0$
    \item $1$
    \item $2$
    \item $3$
\end{enumerate}

\item For a certain reaction, a plot of $\log[R]$ versus time $t$ yields a straight line with a negative slope equal to $-2.303 \times 10^{-3}\text{ s}^{-1}$. The half-life of the reaction is:
\begin{enumerate}
    \item $300.9\text{ s}$
    \item $130.6\text{ s}$
    \item $693\text{ s}$
    \item $100\text{ s}$
\end{enumerate}

\item Which of the following graphs represents a zero-order reaction ($[R]_0$ = initial concentration, $[R]$ = concentration at time $t$, $t_{1/2}$ = half-life)?
\begin{enumerate}
    \item A straight line with positive slope for $[R]$ vs $t$
    \item A straight line passing through the origin for $t_{1/2}$ vs $[R]_0$
    \item A horizontal line parallel to the x-axis for $t_{1/2}$ vs $[R]_0$
    \item A straight line with negative slope for $\ln[R]$ vs $t$
\end{enumerate}

\item A first-order gas-phase reaction $A(g) \rightarrow 2B(g) + C(g)$ starts with pure $A$ at an initial pressure $P_i$. After time $t$, the total pressure of the system is $P_t$. The correct expression for the rate constant $k$ is:
\begin{enumerate}
    \item $k = \frac{2.303}{t} \log \left( \frac{P_i}{2P_i - P_t} \right)$
    \item $k = \frac{2.303}{t} \log \left( \frac{2P_i}{3P_i - P_t} \right)$
    \item $k = \frac{2.303}{t} \log \left( \frac{P_i}{3P_i - P_t} \right)$
    \item $k = \frac{2.303}{t} \log \left( \frac{2P_i}{P_t - P_i} \right)$
\end{enumerate}

\item Consider the thermal decomposition of azoisopropane to hexane and nitrogen at $543\text{ K}$: \\
$( \text{CH}_3 )_2\text{CHN=NCH}(\text{CH}_3)_2(g) \rightarrow \text{N}_2(g) + \text{C}_6\text{H}_{14}(g)$. \\
If the initial pressure is $35.0\text{ mm Hg}$ and after $360\text{ s}$ the total pressure is $54.0\text{ mm Hg}$, the rate constant $k$ is given by:
\begin{enumerate}
    \item $\frac{2.303}{360} \log \left( \frac{35}{19} \right)$
    \item $\frac{2.303}{360} \log \left( \frac{35}{16} \right)$
    \item $\frac{2.303}{360} \log \left( \frac{54}{35} \right)$
    \item $\frac{2.303}{360} \log \left( \frac{35}{54} \right)$
\end{enumerate}

\item Two first-order reactions proceed at $25^\circ\text{C}$ at the same rate. The temperature coefficient of the first reaction is $2$ and that of the second reaction is $3$. The ratio of the rates of the second reaction to the first reaction at $75^\circ\text{C}$ will be:
\begin{enumerate}
    \item $(1.5)^5$
    \item $(1.5)^{50}$
    \item $5 \times 1.5$
    \item $(2/3)^5$
\end{enumerate}

\item A substance decomposes following first-order kinetics. If the fraction remaining after $30\text{ minutes}$ is $\frac{1}{8}$ of the initial amount, what fraction of the initial amount will remain after $50\text{ minutes}$?
\begin{enumerate}
    \item $\frac{1}{16}$
    \item $\frac{1}{32}$
    \item $\frac{1}{64}$
    \item $\frac{1}{24}$
\end{enumerate}

\item For a first-order reaction, if $a$ is the initial concentration of reactant, the concentration of the product formed after time $t$ is given by:
\begin{enumerate}
    \item $a e^{-kt}$
    \item $a(1 - e^{-kt})$
    \item $a(e^{kt} - 1)$
    \item $\frac{a}{1 + kt}$
\end{enumerate}

\item The rate constant of a first-order reaction is $6.93 \times 10^{-3}\text{ min}^{-1}$. Starting with $10\text{ g}$ of the reactant, the rate of reaction after $200\text{ minutes}$ will be:
\begin{enumerate}
    \item $6.93 \times 10^{-3}\text{ g min}^{-1}$
    \item $1.7325 \times 10^{-2}\text{ g min}^{-1}$
    \item $3.465 \times 10^{-2}\text{ g min}^{-1}$
    \item $6.93 \times 10^{-2}\text{ g min}^{-1}$
\end{enumerate}

\item For an $n^{\text{th}}$-order reaction ($n \neq 1$), the half-life $t_{1/2}$ is related to the initial concentration $[R]_0$ as:
\begin{enumerate}
    \item $t_{1/2} \propto [R]_0^{n-1}$
    \item $t_{1/2} \propto [R]_0^{1-n}$
    \item $t_{1/2} \propto [R]_0^n$
    \item $t_{1/2} \propto [R]_0^{-n}$
\end{enumerate}

\item In a zero-order reaction, for every $10^\circ\text{C}$ rise in temperature, the rate is doubled. If the temperature is increased from $10^\circ\text{C}$ to $100^\circ\text{C}$, the rate of the reaction will become:
\begin{enumerate}
    \item $256\text{ times}$
    \item $512\text{ times}$
    \item $64\text{ times}$
    \item $128\text{ times}$
\end{enumerate}

\section*{Section C: Pseudo-First-Order Reactions}

\item During the hydrolysis of $0.01\text{ M}$ ethyl acetate with $10\text{ M}$ water in the presence of dilute $\text{HCl}$:
\[ \text{CH}_3\text{COOC}_2\text{H}_5 + \text{H}_2\text{O} \xrightarrow{\text{H}^+} \text{CH}_3\text{COOH} + \text{C}_2\text{H}_5\text{OH} \]
If the true second-order rate constant is $k = 2.0 \times 10^{-4}\text{ L mol}^{-1}\text{ s}^{-1}$, the pseudo-first-order rate constant $k'$ is:
\begin{enumerate}
    \item $2.0 \times 10^{-3}\text{ s}^{-1}$
    \item $2.0 \times 10^{-6}\text{ s}^{-1}$
    \item $2.0 \times 10^{-2}\text{ s}^{-1}$
    \item $5.0 \times 10^{-3}\text{ s}^{-1}$
\end{enumerate}

\item Which of the following statements is INCORRECT regarding the inversion of cane sugar in aqueous acidic medium?
\begin{enumerate}
    \item The molecularity of the reaction is $2$, while the order of the reaction is $1$.
    \item The concentration of water does not get altered much during the course of the reaction.
    \item The reaction is monitored polarimetrically because sucrose is dextrorotatory while the equimolar mixture of glucose and fructose is laevorotatory.
    \item If the initial concentration of sucrose is doubled, the half-life of the reaction is also doubled.
\end{enumerate}

\item For a second-order reaction $A + B \rightarrow \text{Products}$, the rate law is $\text{Rate} = k[A][B]$. When $[B]_0 \gg [A]_0$, the reaction behaves as a pseudo-first-order reaction. Under this condition, a straight line with a positive slope is obtained when:
\begin{enumerate}
    \item $\ln\left(\frac{[A]_0}{[A]}\right)$ is plotted against $t$, with slope $= k[B]_0$
    \item $\ln\left(\frac{[A]}{[A]_0}\right)$ is plotted against $t$, with slope $= k[B]_0$
    \item $\frac{1}{[A]}$ is plotted against $t$, with slope $= k[B]_0$
    \item $[A]$ is plotted against $t$, with slope $= -k[B]_0$
\end{enumerate}

\section*{Section D: Temperature Dependence, Arrhenius Equation \& Catalyst}

\item According to the Arrhenius equation, $k = A e^{-E_a / RT}$, the slope and y-intercept of the straight-line plot between $\log_{10} k$ (on y-axis) and $\frac{1}{T}$ (on x-axis) are respectively:
\begin{enumerate}
    \item $-\frac{E_a}{R}$ and $\ln A$
    \item $-\frac{E_a}{2.303 R}$ and $\log_{10} A$
    \item $-\frac{2.303 R}{E_a}$ and $\log_{10} A$
    \item $\frac{E_a}{2.303 R}$ and $-\log_{10} A$
\end{enumerate}

\item The rate constant of a reaction increases by $4$ times when the temperature increases from $300\text{ K}$ to $320\text{ K}$. The activation energy ($E_a$) of the reaction is approximately: ($R = 8.314\text{ J K}^{-1}\text{ mol}^{-1}$, $\ln 2 = 0.693$)
\begin{enumerate}
    \item $55.3\text{ kJ mol}^{-1}$
    \item $27.6\text{ kJ mol}^{-1}$
    \item $110.6\text{ kJ mol}^{-1}$
    \item $13.8\text{ kJ mol}^{-1}$
\end{enumerate}

\item In the Arrhenius equation, the term $e^{-E_a / RT}$ represents:
\begin{enumerate}
    \item Total number of collisions per unit volume per second
    \item Fraction of molecules having kinetic energy greater than or equal to activation energy $E_a$
    \item Fraction of collisions with proper steric orientation
    \item Threshold energy of the reactants
\end{enumerate}

\item The rate constants $k_1$ and $k_2$ for two different reactions are $10^{16} e^{-2000/T}$ and $10^{15} e^{-1000/T}$, respectively. The temperature at which $k_1 = k_2$ is:
\begin{enumerate}
    \item $1000\text{ K}$
    \item $\frac{2000}{2.303}\text{ K}$
    \item $\frac{1000}{2.303}\text{ K}$
    \item $2000\text{ K}$
\end{enumerate}

\item For a reversible endothermic reaction $A \rightleftharpoons B$, the enthalpy of reaction is $\Delta H = +40\text{ kJ mol}^{-1}$. If the activation energy of the backward reaction ($E_{a,b}$) is $65\text{ kJ mol}^{-1}$, then the activation energy of the forward reaction ($E_{a,f}$) is:
\begin{enumerate}
    \item $25\text{ kJ mol}^{-1}$
    \item $105\text{ kJ mol}^{-1}$
    \item $40\text{ kJ mol}^{-1}$
    \item $65\text{ kJ mol}^{-1}$
\end{enumerate}

\item A reaction has an activation energy of $57.44\text{ kJ mol}^{-1}$ in the absence of a catalyst. In the presence of a catalyst at $300\text{ K}$, the activation energy is lowered to $45.95\text{ kJ mol}^{-1}$. Assuming the pre-exponential factor $A$ remains unchanged, the rate of the catalyzed reaction compared to the uncatalyzed reaction increases by a factor of: ($R = 8.314\text{ J K}^{-1}\text{ mol}^{-1}$)
\begin{enumerate}
    \item $10$
    \item $100$
    \item $1000$
    \item $20$
\end{enumerate}

\item Which of the following quantities is NOT altered by the addition of a positive catalyst to a reversible chemical reaction?
\begin{enumerate}
    \item Activation energy of the forward reaction
    \item Activation energy of the backward reaction
    \item Rate constants of both forward and backward reactions
    \item Gibbs free energy change ($\Delta G$) and equilibrium constant ($K_{\text{eq}}$)
\end{enumerate}

\item When the temperature of a reacting system is increased by $10\text{ K}$ around room temperature, the reaction rate roughly doubles. The primary reason for this sharp increase in rate is:
\begin{enumerate}
    \item The average kinetic energy of the molecules doubles.
    \item The collision frequency $Z$ doubles.
    \item The fraction of molecules having energy equal to or greater than $E_a$ nearly doubles.
    \item The activation energy $E_a$ decreases to half its value.
\end{enumerate}

\item The decomposition of hydrocarbon follows the equation $k = (4.5 \times 10^{11}\text{ s}^{-1}) e^{-28000\text{ K}/T}$. What is the value of activation energy $E_a$?
\begin{enumerate}
    \item $232.79\text{ kJ mol}^{-1}$
    \item $28.0\text{ kJ mol}^{-1}$
    \item $116.4\text{ kJ mol}^{-1}$
    \item $465.5\text{ kJ mol}^{-1}$
\end{enumerate}

\item At infinite temperature ($T \rightarrow \infty$), the rate constant $k$ of a reaction becomes equal to:
\begin{enumerate}
    \item Zero
    \item $E_a$
    \item $A$ (Arrhenius pre-exponential factor)
    \item $\infty$
\end{enumerate}

\section*{Section E: Collision Theory \& Assertion-Reason}

\item According to the collision theory of bimolecular reactions, the rate of reaction is given by $\text{Rate} = P Z_{AB} e^{-E_a / RT}$. In this expression, the steric factor $P$ accounts for:
\begin{enumerate}
    \item The fraction of molecules possessing threshold energy
    \item The number of collisions per second per unit volume
    \item The proper orientation of colliding molecules at the moment of impact
    \item The temperature dependence of the pre-exponential factor
\end{enumerate}

\item Consider the reaction $\text{CH}_3\text{Br} + \text{OH}^- \rightarrow \text{CH}_3\text{OH} + \text{Br}^-$. According to collision theory, this reaction fails to occur even when colliding molecules possess kinetic energy greater than $E_a$ if:
\begin{enumerate}
    \item $\text{OH}^-$ approaches $\text{CH}_3\text{Br}$ from the side of the $\text{Br}$ atom
    \item $\text{OH}^-$ approaches $\text{CH}_3\text{Br}$ from the side opposite to the $\text{Br}$ atom
    \item The temperature of the reaction mixture is below $298\text{ K}$
    \item No catalyst is added to the reaction vessel
\end{enumerate}

\item What is the threshold energy ($E_{\text{th}}$) of a reaction?
\begin{enumerate}
    \item The minimum extra energy absorbed by reactant molecules to reach the transition state
    \item The minimum total kinetic energy that colliding molecules must possess in order to undergo an effective collision
    \item The difference between the energy of products and the energy of reactants
    \item The average kinetic energy of all reactant molecules at a given temperature
\end{enumerate}

\item \textbf{Assertion (A):} The rate of a zero-order reaction remains constant throughout the course of the reaction until the reactant is completely consumed. \\
\textbf{Reason (R):} For a zero-order reaction, the rate of reaction is independent of the concentration of the reacting species.
\begin{enumerate}
    \item Both (A) and (R) are true and (R) is the correct explanation of (A).
    \item Both (A) and (R) are true but (R) is not the correct explanation of (A).
    \item (A) is true but (R) is false.
    \item Both (A) and (R) are false.
\end{enumerate}

\item \textbf{Assertion (A):} A first-order reaction never reaches $100\%$ completion in finite time. \\
\textbf{Reason (R):} The concentration of the reactant decreases exponentially with time according to $[R] = [R]_0 e^{-kt}$.
\begin{enumerate}
    \item Both (A) and (R) are true and (R) is the correct explanation of (A).
    \item Both (A) and (R) are true but (R) is not the correct explanation of (A).
    \item (A) is true but (R) is false.
    \item Both (A) and (R) are false.
\end{enumerate}

\item \textbf{Assertion (A):} All collisions between reactant molecules with kinetic energy equal to or greater than threshold energy lead to product formation. \\
\textbf{Reason (R):} Effective collisions require both sufficient energy and proper orientation of colliding species.
\begin{enumerate}
    \item Both (A) and (R) are true and (R) is the correct explanation of (A).
    \item Both (A) and (R) are true but (R) is not the correct explanation of (A).
    \item (A) is false but (R) is true.
    \item Both (A) and (R) are false.
\end{enumerate}

\item Match List-I with List-II for a general reaction $R \rightarrow P$:
\begin{center}
\begin{tabular}{|l|l|}
\hline
\textbf{List-I} & \textbf{List-II} \\
\hline
(A) Zero-order reaction & (I) $t_{1/2} \propto [R]_0^{-1}$ \\
(B) First-order reaction & (II) $t_{1/2} \propto [R]_0^0$ \\
(C) Second-order reaction & (III) $t_{1/2} \propto [R]_0^1$ \\
(D) Pseudo-first-order reaction & (IV) $\text{Rate} = k'[R]$ where $k' = k[\text{Excess reactant}]$ \\
\hline
\end{tabular}
\end{center}
Choose the correct answer from the options given below:
\begin{enumerate}
    \item (A)-(III), (B)-(II), (C)-(I), (D)-(IV)
    \item (A)-(II), (B)-(III), (C)-(I), (D)-(IV)
    \item (A)-(III), (B)-(I), (C)-(II), (D)-(IV)
    \item (A)-(I), (B)-(II), (C)-(III), (D)-(IV)
\end{enumerate}

\end{enumerate}

\newpage
\section*{Answer Key \& Brief Hints}

\begin{multicols}{3}
\begin{enumerate}
    \item \textbf{(a)} $-\frac{1}{2}\frac{d[\text{N}_2\text{O}_5]}{dt} = \frac{1}{4}\frac{d[\text{NO}_2]}{dt}$
    \item \textbf{(b)} Rate $= \frac{r_{\text{Br}^-}}{5} = \frac{r_{\text{Br}_2}}{3} = \frac{r_{\text{H}^+}}{6}$
    \item \textbf{(b)} $V \to V/4 \implies [\text{Conc}] \to 4\times; r' = 4 \times 4^2 = 64r$
    \item \textbf{(b)} From I \& IV: $x=1$; from II \& III: $y=2$
    \item \textbf{(b)} $\text{atm}^{1-n}\text{s}^{-1} = \text{atm}^{-3/2}\text{s}^{-1} \implies n = 2.5$
    \item \textbf{(a)} Statement (iii) is false when RDS involves intermediates
    \item \textbf{(b)} $r = k_2[\text{O}_3][\text{O}] = k_2 K_{\text{eq}}[\text{O}_3]^2[\text{O}_2]^{-1}$
    \item \textbf{(a)} $r = k_2[\text{NOCl}_2][\text{NO}] = k_2 K_c[\text{NO}]^2[\text{Cl}_2]$
    \item \textbf{(b)} $2^n = 2\sqrt{2} = 2^{3/2} \implies n = 1.5; 4^{1.5} = 8$
    \item \textbf{(a)} Rate $= \text{bar min}^{-1}$; $k = \text{bar}^{1-1.5}\text{min}^{-1}$
    \item \textbf{(c)} Standard NCERT zero-order example
    \item \textbf{(b)} $k = \frac{0.2-0.15}{10} = 0.005; t_{100\%} = \frac{0.2}{0.005} = 40\text{ min}$
    \item \textbf{(b)} $\frac{t_{80\%}}{t_{40\%}} = \frac{\log(100/20)}{\log(100/60)} = \frac{0.6990}{0.2218} = 3.15 \implies 157.5\text{ min}$
    \item \textbf{(b)} $t_{99.9\%} = \frac{2.303}{k}\log 1000; t_{90\%} = \frac{2.303}{k}\log 10$
    \item \textbf{(c)} $t_{1/2} \propto [R]_0^{1-n} \implies 1-n = -1 \implies n = 2$
    \item \textbf{(a)} $\text{Slope} = -\frac{k}{2.303} \implies k = 5.304 \times 10^{-3}; t_{1/2} = \frac{0.693}{k}$
    \item \textbf{(b)} For zero order, $t_{1/2} = \frac{[R]_0}{2k}$
    \item \textbf{(b)} $P_t = P_i + 2x \implies x = \frac{P_t - P_i}{2}; P_A = P_i - x = \frac{3P_i - P_t}{2}$
    \item \textbf{(b)} $P_t = P_i + x \implies x = 54 - 35 = 19; P_i - x = 16$
    \item \textbf{(a)} $\Delta T = 50^\circ\text{C}$ (5 intervals); $\frac{r_2}{r_1} = \frac{3^5}{2^5} = (1.5)^5$
    \item \textbf{(b)} $\frac{1}{8} = \left(\frac{1}{2}\right)^3 \implies 3t_{1/2} = 30 \implies t_{1/2} = 10\text{ min}$
    \item \textbf{(b)} $[P] = [R]_0 - [R] = a - ae^{-kt} = a(1 - e^{-kt})$
    \item \textbf{(b)} $t_{1/2} = \frac{0.693}{6.93 \times 10^{-3}} = 100\text{ min}; m = \frac{10}{2^2} = 2.5\text{ g}; r = k \cdot m$
    \item \textbf{(b)} Standard integrated rate law relation
    \item \textbf{(b)} $\Delta T = 90^\circ\text{C} \implies 9\text{ intervals} \implies 2^9 = 512$
    \item \textbf{(a)} $k' = k[\text{H}_2\text{O}] = 2.0 \times 10^{-4} \times 10 = 2.0 \times 10^{-3}\text{ s}^{-1}$
    \item \textbf{(d)} Pseudo-first-order $t_{1/2}$ is independent of sucrose conc.
    \item \textbf{(a)} $\ln\left(\frac{[A]_0}{[A]}\right) = (k[B]_0)t$
    \item \textbf{(b)} $\log k = \log A - \frac{E_a}{2.303 R}\left(\frac{1}{T}\right)$
    \item \textbf{(a)} $\ln 4 = \frac{E_a}{8.314}\left(\frac{20}{300 \times 320}\right) \implies E_a \approx 55.3\text{ kJ mol}^{-1}$
    \item \textbf{(b)} Boltzmann factor definition in NCERT
    \item \textbf{(c)} $10^{16}e^{-2000/T} = 10^{15}e^{-1000/T} \implies 10 = e^{1000/T} \implies T = \frac{1000}{2.303}$
    \item \textbf{(b)} $\Delta H = E_{a,f} - E_{a,b} \implies 40 = E_{a,f} - 65 \implies E_{a,f} = 105$
    \item \textbf{(b)} $\log\left(\frac{k_{\text{cat}}}{k_{\text{uncat}}}\right) = \frac{57440 - 45950}{2.303 \times 8.314 \times 300} = 2 \implies \frac{k_{\text{cat}}}{k_{\text{uncat}}} = 100$
    \item \textbf{(d)} Catalysts do not alter thermodynamic state functions or $K_{\text{eq}}$
    \item \textbf{(c)} Area under Maxwell-Boltzmann curve beyond $E_a$ doubles
    \item \textbf{(a)} $\frac{E_a}{R} = 28000\text{ K} \implies E_a = 28000 \times 8.314 = 232.79\text{ kJ mol}^{-1}$
    \item \textbf{(c)} As $T \to \infty$, $e^{-E_a/RT} \to e^0 = 1 \implies k \to A$
    \item \textbf{(c)} $P$ is the orientation/probability factor
    \item \textbf{(a)} Front-side attack causes electrostatic repulsion with $\text{Br}$
    \item \textbf{(b)} $E_{\text{th}} = E_{\text{average}} + E_a$
    \item \textbf{(a)} $\text{Rate} = k[R]^0 = k$ (constant)
    \item \textbf{(a)} $[R] = 0$ only when $t \to \infty$
    \item \textbf{(c)} Assertion is false because orientation is also mandatory
    \item \textbf{(a)} Direct match of half-life dependencies
\end{enumerate}
\end{multicols}

\end{document}